%% BLOQUE: datos
Se analizan 390 pacientes con 17 variables: colesterol (\texttt{chol}), glucosa
estabilizada (\texttt{stab.glu}), HDL, razón colesterol/HDL (\texttt{ratio}), edad,
sexo, localidad, talla, peso, complexión (\texttt{frame}), dos mediciones de presión
arterial, cintura, cadera y tiempo desde la última comida (\texttt{time.ppn}). Sólo el
15.4\,\% son diabéticos (60 casos): en prueba hay 15 diabéticos, por lo que cada
diabético mal clasificado cambia el recall en 6.7 puntos, y en cada pliegue de
validación hay unos 9, lo que explica desviaciones de 0.06 a 0.12 en el F1 macro. La
glucosa estabilizada separa mucho mejor las clases que cualquier variable de
Potabilidad: su mediana es 186\,mg/dL en diabéticos y 86.5\,mg/dL en no diabéticos.
Tras la codificación one-hot hay 21 atributos; la información mutua de la glucosa
(0.228) es tres veces la del siguiente, \texttt{ratio} (0.063).

%% BLOQUE: KNN
Con todos los atributos, KNN es la familia más débil (F1 macro 0.719 en el estándar).
Aumentar a $k=15$ empeora (0.681): con sólo unos 36 diabéticos en cada ajuste, en un
vecindario de 15 suele haber mayoría de no diabéticos y el voto favorece la clase
mayoritaria. La distancia Manhattan también empeora (0.636) y ponderar por distancia dio
exactamente las mismas métricas que el estándar en los cinco pliegues. El cambio que sí
importa es la selección de atributos: con los 11 de mayor información mutua el F1 macro
sube a 0.797 (+0.078), la mayor mejora de este dataset. Al eliminar variables poco
informativas como la talla, la localidad, el sexo o las segundas mediciones de presión
(ausentes en el 65\,\% de los pacientes e imputadas con la mediana), la distancia pasa
a reflejar sobre todo la glucosa. En prueba detecta 7 de 15 diabéticos (recall 0.467)
con precisión 0.70. Su AUC (0.754) es la más baja porque, con 5 vecinos, la
probabilidad sólo puede tomar seis valores (0, 0.2, \dots, 1), lo que produce una curva
ROC escalonada.

%% BLOQUE: DT
Exigir al menos 10 observaciones por hoja es de nuevo la mejor variante (0.806, +0.024),
igual que en Potabilidad. Limitar la profundidad a 5 da un resultado parecido (0.799)
pero con una variación entre pliegues mucho mayor (0.117 frente a 0.070). Ponderar las clases empeora (0.751): con
tan pocos diabéticos, darles más peso hace que el árbol cree hojas para casos
aislados. La selección de atributos mejora un poco (0.816) y es la versión final. El
árbol resultante es pequeño (profundidad 5, 11 hojas) y fácil de leer: su primera regla
es \texttt{stab.glu} $>158$\,mg/dL $\Rightarrow$ diabetes, y en el rango de 107.5 a
158\,mg/dL marca como diabéticos a los pacientes con colesterol mayor que 243.5\,mg/dL. La importancia por
permutación confirma que en prueba el árbol depende exclusivamente de la glucosa (caída
de 0.252 en F1 macro; el resto, cero). En prueba detecta 8 de 15 diabéticos (recall
0.533, precisión 0.727) con exactitud 89.8\,\% y AUC 0.907. El umbral de 158\,mg/dL lo
aprendió el árbol de estos datos; no es un criterio diagnóstico. Como referencia, el
criterio clínico para una glucosa casual es 200\,mg/dL con síntomas \cite{ada}.

%% BLOQUE: RF
En validación cruzada, Random Forest es el mejor algoritmo del dataset. Su
configuración estándar (0.836) supera a las tres variantes: 500 árboles no mejora
(0.818), la ponderación de clases tampoco (0.816) y exigir 5 observaciones por hoja
empeora claramente (0.763), porque impide aislar a los pocos diabéticos en hojas
propias. Con selección de atributos alcanza 0.840 y una desviación menor (0.066 frente
a 0.087), así que esa es la versión final. En prueba tiene exactitud 88.8\,\%, F1
macro 0.764 y AUC 0.892; detecta 8 de 15 diabéticos y tiene un falso positivo más que
DT, XGBoost y AdaBoost (4 en lugar de 3). Por eso su F1 de prueba es ligeramente menor,
aunque fue el elegido por validación. Esa diferencia de un paciente está muy por debajo
de la incertidumbre del experimento. La importancia interna muestra que, además de la
glucosa (0.454), el bosque reparte peso entre \texttt{ratio}, HDL y edad, pero en
prueba sólo la glucosa produce una caída grande al permutarse.

%% BLOQUE: XGBoost
Bajar la tasa de aprendizaje a 0.05 es la mejor variante (0.821, +0.029): con sólo 292
pacientes de entrenamiento, que cada árbol aporte menos funciona como regularización.
Reducir la profundidad a 3 ayuda algo menos (0.804) y ponderar las clases no ayuda
(0.789). A diferencia de las otras familias, la selección de atributos empeora a XGBoost
(0.803): el algoritmo ya elige en cada división la variable más útil y quitarle
variables secundarias le resta poco ruido y algo de información. En prueba obtiene la
mejor AUC del estudio (0.929), exactitud 89.8\,\% y F1 macro 0.778, con la misma matriz
de confusión que DT y AdaBoost. Su F1 de entrenamiento (0.993) indica que aun con tasa
baja casi memoriza el entrenamiento.

%% BLOQUE: AdaBoost
Las variantes de AdaBoost quedan todas entre 0.790 y 0.799, sin diferencias claras. La
tasa 0.5 es la de mayor media (0.799), pero también la de mayor desviación (0.122). La
ponderación de clases logra el mayor recall macro de validación (0.801), aunque no el
mayor F1. Con selección de atributos sube a 0.803 y esa es la versión final. En prueba
coincide con DT y XGBoost (exactitud 89.8\,\%, F1 macro 0.778, 8 de 15 diabéticos) y
tiene AUC 0.922. Es el ensamble con menor sobreajuste (F1 0.887 en entrenamiento frente
a 0.778 en prueba). En este dataset los stumps sí funcionan, a diferencia de
Potabilidad: un único corte en la glucosa ya separa bien las clases, y los siguientes
stumps afinan con el colesterol (importancia interna 0.120) y el peso (0.106).

%% BLOQUE: sintesis
Los cuatro métodos basados en árboles quedan entre 0.80 y 0.84 de F1 macro en
validación cruzada, con desviaciones de 0.07 a 0.12. Las diferencias entre ellos son
menores que esa variación, así que el orden no es concluyente; RF se eligió por tener
la mayor media. En prueba todos superan claramente la referencia (F1 macro 0.459) y DT,
XGBoost y AdaBoost producen la misma matriz de confusión: 8 diabéticos detectados, 7 no
detectados y 3 falsos positivos.

Los 7 diabéticos que no detectan DT, RF, XGBoost y AdaBoost son exactamente los mismos
(KNN falla esos 7 y uno más), y todos tienen glucosa estabilizada entre 90 y 130\,mg/dL, mientras que los 8 detectados
tienen 196\,mg/dL o más. La etiqueta se definió con la hemoglobina glicosilada, que
refleja la glucosa promedio de los últimos dos o tres meses, y puede estar elevada
aunque una medición puntual de glucosa sea normal. Como \texttt{glyhb} se excluyó para
evitar fuga, esos pacientes no se distinguen de los no diabéticos con las variables
disponibles. Es un límite de la información, no de los algoritmos: ningún ajuste de
hiperparámetros los recuperaría sin aumentar mucho los falsos positivos.

Las AUC de 0.89 a 0.93 indican que los modelos ordenan bien a los pacientes por riesgo.
Bajar el umbral de decisión de 0.5 aumentaría el recall de diabetes a cambio de más
falsos positivos. Ese umbral debería elegirse con validación cruzada y no con el
conjunto de prueba, por lo que se deja como trabajo futuro. A diferencia de
Potabilidad, aquí la selección de atributos ayudó a cuatro de cinco familias
(Figura~\ref{fig:seleccion-diabetes}): la información relevante se concentra en una
variable con relación monótona con la clase, y la información mutua la detecta sin
problema.
